\documentclass[10pt,a4paper]{amsart}
\usepackage[margin=19mm]{geometry}
\usepackage{amsmath,amssymb,mathtools}
\usepackage[scr=rsfs,cal=euler]{mathalfa}
\usepackage{xfrac}
\usepackage[unicode,hidelinks,bookmarks=false]{hyperref}
\newcommand{\ie}{i.e.\ }
\newcommand{\cf}{cf.\ }
\newcommand{\resp}{respectively\ }
\newcommand{\Subsection}{\subsection}
\providecommand{\nobreakdash}{\nobreak\hbox{-}}
\setlength{\emergencystretch}{2em}
\multlinegap0pt
\usepackage{amssymb}
\usepackage{tensor}
\usepackage{mathtools}
\usepackage{tikz}
\usepackage{bm}
\usepackage{dsfont}
\newtheorem{lemma}{Lemma}
\usepackage{cleveref}
\crefname{lemma}{Lemma}{Lemmas}
\DeclareMathOperator{\Tr}{Tr}
\newcommand{\boundarymesh}[0]{\mathcal{U}}
\DeclareMathOperator{\curl}{curl}
\newcommand{\mesh}[0]{\mathcal{T}}
\DeclareMathOperator{\notsubsimplex}{\ntrianglelefteq}
\DeclareMathOperator{\subsimplex}{\unlhd}
\newcommand{\unitvec}[1]{\hat{#1}}
\title{full title}
\author{Pablo D. Brubeck, Yizhou Liang, Charles Parker}
\date{Source: arXiv:2605.31348v1 (CC BY 4.0)}
\begin{document}
\maketitle

\noindent\emph{Source and licence.} full title. Version arXiv:2605.31348v1 (CC BY 4.0), distributed by arXiv under the Creative Commons Attribution 4.0 International licence, http://creativecommons.org/licenses/by/4.0/, as recorded in the arXiv licence metadata for this exact version and shown on the abstract page https://arxiv.org/abs/2605.31348v1. Full licence notice: \texttt{licenses/arxiv-2605.31348-CC-BY-4.0.txt}.\par\smallskip

\noindent\textit{Editorial scope.} Counted unit: the article's lemma lemma (source0.tex lines 1749--1753). The article's complete original proof of it is delivered with it (source0.tex lines 1754--1843, one \texttt{proof} environment of 90 lines). No mathematics is written, repaired or re-derived: the statement and the proof are the article's own lines, in the article's own notation, and no retained line is substituted or re-flowed.  Retained uncounted as the source-local context the proof needs: lines 1545--1568.  Nothing was omitted from any retained span, so the package carries no omission record.  No retained line contains a citation, so the wrapper carries no bibliography and no bibliography entry.  No retained line contains a figure environment, an image-inclusion command or drawing code, so no image is delivered and the package declares no asset dependency.  Editorial formatting only, in addition: an AMS \texttt{amsart} preamble that restates the article's own declarations and macros used by the retained lines, namely the environments \texttt{lemma} on separate counters, together with the standard packages those lines need.  The retained environments print in this document as Lemma 1 in order of appearance, whereas the article prints the same items with the numbering of its own full document.  The complete native source of the article as distributed in the arXiv e-print for this exact version is bundled unchanged as \texttt{sources/201714/source0.tex}.
\begin{lemma}
	\label{lem:support-skeleton-extended-bubble}
	Let $\tau \in \Delta(\mesh)$. \\
	\noindent (i) For $v \in \mathbb{B}^k(\tau)$,
	there holds
	\begin{align}
		\label{eq:modified-bubble-support}
		\Tr^k_{\eta} v = 0 \qquad \forall \eta \in 
		\left( \bigcup_{\ell=0}^{\dim \tau} \Delta_{\ell}(\mesh) 
		\cup \bigcup_{\ell=\dim \tau + 1}^{3} 
		\{ \rho \in \Delta_{\ell}(\mesh) : \tau \notsubsimplex \rho \}
		\right) \setminus \{\tau\}.
	\end{align}
	
	\noindent (ii)  For $v \in \mathcal{S}^k$, 
	there holds
	\begin{align}
		\label{eq:skeletal-support}
		\mathcal{I}^k_{\eta}(v) = 0 \qquad \forall \eta \in \Delta_k(\tau)
		\implies \Tr^k_{\tau} v = 0.
	\end{align}
	In particular, $\Tr^k_{\tau} v = 0$ if
	$\tau \in \bigcup_{\ell = 0}^{k-1} \Delta_{\ell}(\mesh)$.
\end{lemma}

\begin{lemma}
	For $k \in 0:3$, $\mathcal{S}^k_{\Gamma_0} \subset V_{\Gamma_0}^{k, h}$
	and $\mathbb{B}_{\Gamma_0}^k(\tau) \subset V_{\Gamma_0}^{k, h}$ for all 
	$\tau \in \Delta(\mesh)$.
\end{lemma}

\begin{proof}
	Let $v \in \mathcal{S}^k_{\Gamma_0}$. For any $f \in \boundarymesh$,
	$\mathcal{I}^k_{\eta}(v) = 0$ for all $\eta \in \Delta_k(f)$, and so 
	$\Tr^k_{f} v = 0$ thanks to 
	\cref{lem:support-skeleton-extended-bubble}. 
	Thus, $v \in V_{\Gamma_0}^{k, h}$.
	
	Now let $\tau \in \Delta(\mesh)$ and 
	$v \in \mathbb{B}_{\Gamma_0}^k(\tau)$. 
	If $\tau \notin \Delta(\boundarymesh)$,
	then $\Tr^k_{f} v = 0$ for all $f \in \boundarymesh$
	thanks to \cref{lem:support-skeleton-extended-bubble},
	and so $v \in V_{\Gamma_0}^{k, h}$. Now suppose that 
	$\tau \in \Delta_0^{\flat}(\boundarymesh)$. Applying 
	\cref{lem:support-skeleton-extended-bubble} once again gives 
	$\Tr_{f}^k v = 0$ for $f \in \boundarymesh$
	with $\tau \notsubsimplex f$. Consequently, the final case to verify 
	is $f \in \boundarymesh$ with $\tau \subsimplex f$.
	
	For $k = 0$, expanding the definition of the harmonic inner products
	in the condition $v \in \mathbb{B}_{\Gamma_0}^k(\tau)$ shows that	
	$v_f := v|_{f} \in \mathcal{P}_{p+2}(f)$ satisfies 
	\begin{subequations}
		\label{eq:proof:Pp2-face-dofs}
		\begin{alignat}{2}
			D_f^{\alpha} v_f (z) &= 0 \qquad & &\forall |\alpha| \leq 2, 
			\ \forall z \in \Delta_0(f), \\
			(\partial_{\unitvec{t}_{\partial f}}v_f, \partial_{\unitvec{t}_{\partial f}}w)_{L^2(e)} &= 0 
			\qquad & &\forall w \in \mathcal{P}_{p+2}(e) \cap H^3_0(e), \ 
			\forall e \in \Delta_1(f), \\
			(\partial_{\unitvec{n}_{\partial f}} v_{f}, w )_{L^2(e)} &= 0 
			\qquad & &\forall w \in \mathcal{P}_{p+1}(e) \cap H^2_0(e), \ 
			\forall e \in \Delta_1(f), \\ 
			(\operatorname{grad}_fv_f, \operatorname{grad}_fw)_{L^2(f)} &=0 
			\qquad & &\forall w \in \mathcal{P}_{p+2}(f) \cap H^2_0(f),
		\end{alignat}	
	\end{subequations}
	where we recall $D_f$ denotes the surface differential.
	Since the above degrees of freedom are unisolvent on $\mathcal{P}_{p+2}(f)$,
	$v_{f} \equiv 0$. Similarly, 
	$u_f = \partial_{n_f} v|_f \in \mathcal{P}_{p+1}(f)$ satisfies
	\begin{subequations}
		\label{eq:proof:Pp1-face-dofs}
		\begin{alignat}{2}
			D_f^{\alpha} u_f (z) &= 0 \qquad & &\forall |\alpha| \leq 1, 
			\ \forall z \in \Delta_0(f), \\
			(u_f, w)_{L^2(e)} &= 0 
			\qquad & &\forall w \in \mathcal{P}_{p+1}(e) \cap H^2_0(e), \ 
			\forall e \in \Delta_1(f),  \\
			\label{eq:proof:Pp1-face-dofs-interior-face}
			(u_f, w)_{L^2(f)} &=0 
			\qquad & &\forall w \in \mathcal{P}_{p+1}(f) \cap H^1_0(f),
		\end{alignat}	
	\end{subequations}	
	and so $u_f \equiv 0$. As a result, $v \in V_{\Gamma_0}^{0, h}$.
	
	For $k=1$, each component of 
	$v_f := v|_{f} \in \mathcal{P}_{p+1}(f)^3$ also satisfies 
	\begin{subequations}
		\label{eq:proof:Pp3-face-dofs}
		\begin{alignat}{2}
			D_f^{\alpha} v_f (z) &= 0 \qquad & &\forall |\alpha| \leq 1, 
			\ \forall z \in \Delta_0(f), \\
			(v_f, w)_{L^2(e)} &= 0 
			\qquad & &\forall w \in \mathcal{P}_{p+1}(e)^3 \cap H^2_0(e)^3, \ 
			\forall e \in \Delta_1(f),  \\
			\label{eq:proof:Pp3-face-dofs-interior-face}
			(v_f, \operatorname{grad}_f w)_{L^2(f)} &=0 
			\qquad & &\forall w \in \mathcal{P}_{p+2}(f) \cap H^2_0(f),\\
			(\operatorname{rot}_fv_f, \operatorname{rot}_fw)_{L^2(f)} &=0 
			\qquad & &\forall w \in \mathcal{P}_{p+1}(f)^3 \cap H^1_0(f)^3,
		\end{alignat}	
	\end{subequations}	 and so 
	$v_f \equiv 0$. Moreover, $u_f := \curl v|_f
	\in \mathcal{P}_p(f)^3$ satisfies	
	\begin{subequations}
		\label{eq:proof:Pp-face-dofs}
		\begin{alignat}{2}
			u_f (z) &= 0 
			\qquad & &\forall z \in \Delta_0(f), \\
			(u_f, w)_{L^2(e)} &= 0 
			\qquad & &\forall w \in \mathcal{P}_{p}(e)^3 
			\cap H^1_0(e)^3, \ \forall e \in \Delta_1(f), \\
			(u_f, w)_{L^2(f)} &=0 
			\qquad & &\forall w \in \mathcal{P}_{p}(f)^3 
			\cap H^1_0(f)^3,
		\end{alignat}	
	\end{subequations}
	and so $u_f \equiv 0$. Consequently, $v \in V_{\Gamma_0}^{1, h}$.
\end{proof}

\end{document}
